\chapter{Palavras e conceitos}

\begin{objetivo}
Conhecer as palavras hebraicas e gregas do tema, quantas vezes e onde
aparecem, e anotar ideias do tema que aparecem sem essas palavras.
\end{objetivo}

\section{Vocabulário}
\vocabulariodotema

Somadas, \hb{דָּם} e \gr{αἷμα} aparecem \temaNOcorrencias\ vezes (lista
completa no apêndice). Levítico concentra 88 das 361 ocorrências de \hb{דָּם};
Hebreus, 21 das 97 de \gr{αἷμα}: são os dois livros que mais tratam de
sacrifícios.

\begin{alerta}[Palavra não é conceito]
Nem toda ocorrência de \enquote{sangue} trata do mesmo aspecto: em
\rb{Êx 7:17–21} é a água do Nilo transformada; em muitos textos dos profetas,
a culpa por homicídios. E há textos do tema sem a palavra, por exemplo sobre
a vida como dom de Deus.
\end{alerta}

\campovazio{Sinônimos e expressões}{ex.: “derramar”, “vida”, “altar”, “expiação”}
\campovazio{O tema sem a palavra}{}
